\documentclass[10pt,a4paper]{article}
\usepackage[T1]{fontenc}
\usepackage[utf8]{inputenc}
\usepackage[brazilian]{babel}
\usepackage[a4paper,margin=1.5cm]{geometry}
\usepackage{booktabs}
\usepackage{xcolor}
\usepackage{microtype}
\usepackage[hidelinks]{hyperref}
\pagestyle{empty}
\setlength{\parindent}{0pt}
\setlength{\parskip}{2pt}
% Paleta monocromática: só a escala gray do xcolor, para impressão em toner.
\definecolor{cinza}{gray}{0.35}
\definecolor{cinzaclaro}{gray}{0.94}
\definecolor{borda}{gray}{0.75}
\definecolor{fundocard}{gray}{0.98}
\definecolor{barras}{gray}{0.15}
\newcommand{\barra}[1]{\textcolor{barras}{\rule{#1mm}{2.4pt}}}
% Cartão de indicador: rótulo, número em destaque e rodapé de contexto.
\newcommand{\kpicard}[3]{%
  \fcolorbox{borda}{fundocard}{%
    \begin{minipage}[c][1.4cm][c]{\dimexpr0.32\textwidth-2\fboxsep-2\fboxrule\relax}
      \centering
      {\scriptsize\bfseries\color{cinza} #1}\\[2pt]
      {\Large\bfseries #2}\\[1.5pt]
      {\tiny\color{cinza} #3}%
    \end{minipage}%
  }%
}
\begin{document}

{\Large\bfseries Tapiraí --- climatologia das médias}\\[2pt]
{\footnotesize\color{cinza} Estação CIIAGRO/IAC ciiagro:82 \,\textbullet\, 231 meses (2005-08 a 2026-10)}\\[1pt]
{\footnotesize \textbf{15 anos com os 12 meses fechados} --- médias do período medido, não previsão do tempo.}

\vspace{4pt}
\colorbox{cinzaclaro}{%
  \parbox{\dimexpr\textwidth-2\fboxsep}{%
    \footnotesize
    \textbf{Leia antes de usar:} os números abaixo são \textbf{médias de vários anos}, obtidas das medições mensais do CIIAGRO/IAC. \textbf{Não é previsão meteorológica} para os próximos dias nem para esta safra.\par\vspace{2pt}
    \textbf{Aplicação prática:} este é o comportamento médio esperado da estação ---
    base técnica para o calendário agrícola, janelas de calagem, o preparo do solo e
    o manejo conservacionista.%
  }%
}
\vspace{4pt}
\noindent
\kpicard{Chuva anual média}{1.592 mm/ano}{em todo o ano}\hfill
\kpicard{Mês mais chuvoso}{222 mm}{em janeiro}\hfill
\kpicard{Mês mais seco}{65 mm}{em agosto}

\vspace{2pt}
\noindent
\kpicard{Temperatura média}{18,4 \textdegree{}C}{média dos 12 meses}\hfill
\kpicard{Mês mais quente}{22,9 \textdegree{}C}{em fevereiro}\hfill
\kpicard{Mês mais frio}{14,0 \textdegree{}C}{em julho}

\vspace{4pt}
\begin{center}
{\bfseries\small Tabela 1: Chuva média por mês do ano}\\[2pt]
{\small
\begin{tabular}{@{}l r l r r r@{}}
\toprule
mês & chuva (mm) & perfil & t mín & t méd & t máx \\
\midrule
janeiro    &   222 & \barra{26.0}   & 18,0 & 22,8 & 27,6 \\
fevereiro  &   185 & \barra{21.6}   & 17,9 & 22,9 & 28,0 \\
março      &   157 & \barra{18.4}   & 16,7 & 21,4 & 26,1 \\
abril      &   100 & \barra{11.7}   & 14,8 & 19,5 & 24,2 \\
maio       &    81 & \barra{9.5}    & 11,6 & 16,7 & 21,9 \\
junho      &    81 & \barra{9.5}    &  9,6 & 15,2 & 20,7 \\
julho      &    80 & \barra{9.4}    &  7,3 & 14,0 & 20,8 \\
agosto     &    65 & \barra{7.6}    &  9,6 & 14,1 & 21,2 \\
setembro   &   118 & \barra{13.8}   & 11,6 & 15,6 & 22,1 \\
outubro    &   135 & \barra{15.9}   & 13,3 & 17,4 & 21,6 \\
novembro   &   170 & \barra{19.9}   & 15,3 & 19,8 & 24,4 \\
dezembro   &   185 & \barra{21.7}   & 17,1 & 21,8 & 26,5 \\
\midrule
ano & 1.592 & & & 18,4 & \\
\bottomrule
\end{tabular}}
\end{center}

\begin{center}
{\bfseries\small Tabela 2: Chuva por estação do ano}\\[2pt]
{\small
\begin{tabular}{@{}l l r r r l l@{}}
\toprule
estação & meses & total (mm) & \% do ano & média/mês & maior mês & menor mês \\
\midrule
Verão      & dez--fev &   591 &  37\% &  197 & 222 (jan) & 185 (fev) \\
Outono     & mar--mai &   337 &  21\% &  112 & 157 (mar) & 81 (mai) \\
Inverno    & jun--ago &   226 &  14\% &   75 & 81 (jun) & 65 (ago) \\
Primavera  & set--nov &   423 &  27\% &  141 & 170 (nov) & 118 (set) \\
\bottomrule
\end{tabular}}
\end{center}

\vspace{1pt}
{\footnotesize\textbf{Antes dos solstícios} (acumulado desde janeiro):
inverno: janeiro a maio somam 743 mm (47\% do ano) \,\textbullet\, verão: junho a novembro somam 649 mm (41\% do ano). Dezembro fica depois do solstício de verão e entra só no total
anual.}

\vspace{2pt}
{\bfseries\small Como usar no planejamento nesta estação}

\footnotesize\begin{itemize}\setlength{\itemsep}{0.5pt}
  \item \textbf{Amostragem de solo e preparo mecânico}: Em abril, maio, junho, julho e agosto a chuva média fica entre 65 e 100 mm/mês \textemdash{} é a janela de menor encharcamento para retirar terra para laudo, arar e incorporar calcário com antecedência. Como até o mês mais seco passa de 60 mm, o solo raramente seca: a janela vale para entrar com máquina, não para preparo em solo seco.
  \item \textbf{Calagem em cobertura (pasto e perene já formado)}: A partir de outubro, quando a média volta a passar de 120 mm/mês: aí há água para o corretivo descer no perfil sem enxurrada.
  \item \textbf{Plantio e transplante}: O verão (dezembro a fevereiro) concentra 591 mm, 37\% da chuva do ano \textemdash{} favorece o pegamento de mudas sem irrigação pesada; em áreas declivosas, reforce a conservação de solo nesse período.
  \item \textbf{Excesso de chuva e conservação}: janeiro passa de 200 mm em média (222 mm): canteiros altos, drenagem limpa e cobertura morta antes da temporada, para não perder solo e adubo no escoamento.
  \item \textbf{Calor, frio e fitossanidade}: fevereiro tem a maior média (22,9 \textdegree{}C) e julho a menor (14,0 \textdegree{}C, com mínima média de 7,3 \textdegree{}C). Antecipe o monitoramento de doenças fúngicas nos meses quentes e úmidos e proteja culturas sensíveis ao frio nos mais frios.
\end{itemize}

\vspace{1pt}
{\bfseries\small Vale para qualquer município da região}

\scriptsize\begin{itemize}\setlength{\itemsep}{0.5pt}
  \item \textbf{Média não é previsão}: Descrevem o padrão de vários anos, não o tempo deste ano: use para escolher a janela de cada operação, e decida no campo.
  \item \textbf{Corrija antes de precisar}: Calcário reage em meses; aplicado na janela seca, chega ao plantio já reagido.
  \item \textbf{Chuva em mm não é água no solo}: Textura, declive e cobertura morta decidem quanto do que cai chega à raiz.
  \item \textbf{Série curta é pista, não padrão}: Poucos anos fechados oscilam muito: compare com a estação vizinha antes de fixar calendário.
\end{itemize}

\vspace{1pt}
{\tiny\color{cinza}
Fonte: CIIAGRO/IAC (\href{http://ciiagro.org.br}{ciiagro.org.br}), série mensal;
chuva é o total do mês. Painel:
\href{https://climalitoralsul.netlify.app/\#tapirai}%
{\texttt{climalitoralsul.netlify.app/\#tapirai}}.
Umidade relativa omitida: a fonte usa \texttt{0} para ``sem sensor''.  3 leitura(s) fora da faixa de sanidade excluída(s) da média (tmin 2, chuva 1); o banco as preserva.
médias do período 2005\textendash{}2026 (19.2 anos de dado).
Gerado a partir de \texttt{dados.json} de 2026-10-30.}

\end{document}
